\section{Математическое описание}

\subsection{Конечный автомат}

Конечным автоматом называется математическая модель дискретного устройства,
которое в каждый момент времени находится в одном из конечного числа состояний.
Детерминированный конечный автомат задаётся набором
\begin{equation}
    A = (S, \Sigma, Y, s_0, \delta, \lambda),
\end{equation}
где $S$ --- конечное множество состояний, $\Sigma$ --- множество входных
сигналов, $Y$ --- множество выходных сигналов, $s_0 \in S$ --- начальное
состояние, $\delta \colon S \times \Sigma \to S$ --- функция переходов,
$\lambda \colon S \times \Sigma \to Y$ --- функция выходов.

Клетку рассматриваемого клеточного автомата можно интерпретировать как автомат
Мура: её выходом является текущее состояние, а входом --- состояния соседних
клеток. Состояние на следующем такте однозначно определяется локальной функцией
перехода.

\subsection{Клеточный автомат}

Клеточный автомат представляет собой регулярную решётку одинаковых конечных
автоматов, функционирующих синхронно в дискретном времени. Состояние каждой
клетки на следующем такте зависит от её текущего состояния и состояний клеток
заданной окрестности. Рассматриваемый автомат задаётся четвёркой
\begin{equation}
    \mathcal{A} = (L, S, \mathcal{N}, f),
\end{equation}
где $L$ --- клеточное пространство, $S$ --- множество состояний клетки,
$\mathcal{N}$ --- окрестность клетки, $f$ --- локальная функция перехода.

Поле представляет собой квадратную решётку размера $n \times n$:
\begin{equation}
    L = \mathbb{Z}_n \times \mathbb{Z}_n,
    \qquad
    \mathbb{Z}_n = \{0, 1, \ldots, n-1\}.
\end{equation}
Множество состояний клетки является бинарным:
\begin{equation}
    S = \{0,1\},
\end{equation}
где состояние $0$ соответствует мёртвой клетке, а состояние $1$ --- живой.
Состояние клетки с координатами $(i,j)$ в момент времени $t$ обозначается
$s^t_{i,j} \in S$, а состояние всего поля --- матрицей
\begin{equation}
    C^t = \left(s^t_{i,j}\right)_{i,j=0}^{n-1}.
\end{equation}

\subsection{Окрестность фон Неймана}

Для каждой клетки используется окрестность фон Неймана, содержащая саму
клетку и четыре клетки, имеющие с ней общую сторону. Для однозначного задания
аргументов функции перехода вводится упорядоченная окрестность
\begin{equation}
    \mathcal{N}_{i,j}^{t} =
    \left(
        s^t_{i,j},
        s^t_{i-1,j},
        s^t_{i+1,j},
        s^t_{i,j-1},
        s^t_{i,j+1}
    \right).
    \label{eq:neumann-neighborhood}
\end{equation}
Компоненты кортежа~\eqref{eq:neumann-neighborhood} записаны в порядке
«центр, верх, низ, лево, право» в соответствии с лекцией.

\subsection{Тороидальные граничные условия}

На поле заданы тороидальные граничные условия. Арифметика индексов выполняется
по модулю $n$:
\begin{align}
    i_{\mathrm{up}} &= (i-1) \bmod n,
    & i_{\mathrm{down}} &= (i+1) \bmod n, \\
    j_{\mathrm{left}} &= (j-1) \bmod n,
    & j_{\mathrm{right}} &= (j+1) \bmod n.
\end{align}
Следовательно, верхняя и нижняя, а также левая и правая границы поля попарно
отождествляются, и каждая клетка имеет одинаковое число соседей.

У клетки верхней строки верхний сосед находится в последней строке того же
столбца. Аналогично левый сосед клетки первого столбца находится в последнем
столбце той же строки. В частности, для левой верхней клетки справедливо
\begin{equation}
    s^{t+1}_{0,0} = f\left(
        s^t_{0,0},
        s^t_{n-1,0},
        s^t_{1,0},
        s^t_{0,n-1},
        s^t_{0,1}
    \right).
\end{equation}

\subsection{Правило перехода}

Локальная функция перехода имеет вид
\begin{equation}
    f \colon S^5 \to S.
\end{equation}
Поскольку окрестность состоит из пяти бинарных клеток, существует
$2^5 = 32$ возможных набора её состояний. Каждому элементу множества $S^5$
ставится в соответствие номер строки таблицы истинности
\begin{equation}
    k^t_{i,j} =
        16s^t_{i,j}
        + 8s^t_{i-1,j}
        + 4s^t_{i+1,j}
        + 2s^t_{i,j-1}
        + s^t_{i,j+1},
    \qquad 0 \leq k^t_{i,j} \leq 31,
    \label{eq:combination-index}
\end{equation}
То есть состояния клеток окрестности рассматриваются как разряды одного
пятибитного двоичного числа.

Правило перехода кодируется 32-битным числом
\begin{equation}
    N = 37 \cdot 11 \cdot 2006 \cdot 17 \cdot 4 = 55\,518\,056.
\end{equation}
Его двоичная запись, дополненная ведущими нулями до 32 разрядов, равна
\begin{equation}
    N = 00000011010011110010001101101000_2.
\end{equation}
Значение функции для строки с номером $k^t_{i,j}$ равно соответствующему биту
числа $N$:
\begin{equation}
    f\left(
        s^t_{i,j},s^t_{i-1,j},s^t_{i+1,j},s^t_{i,j-1},s^t_{i,j+1}
    \right)
    = b_{k^t_{i,j}}
    = \left\lfloor \frac{N}{2^{k^t_{i,j}}} \right\rfloor \bmod 2.
    \label{eq:transition-rule}
\end{equation}
В обычной записи двоичного числа слева расположен бит $b_{31}$, а справа ---
бит $b_0$. Поэтому при заполнении таблицы от строки $k=0$ к строке $k=31$
биты двоичной записи $N$ читаются справа налево.

Полная таблица истинности функции перехода приведена в
табл.~\ref{tab:truth-table}.
\begin{table}[p]
    \centering
    \caption{Таблица истинности функции перехода}
    \label{tab:truth-table}
    \small
    \begin{tblr}{
        colspec = {Q[c] Q[c] Q[c] Q[c] Q[c] Q[c] Q[c]},
        row{1} = {font=\bfseries},
        rowsep = 0pt,
        hlines,
        vlines,
    }
        $k$ & $s^t_{i,j}$ & $s^t_{i-1,j}$ & $s^t_{i+1,j}$ &
        $s^t_{i,j-1}$ & $s^t_{i,j+1}$ & $f$ \\
         0 & 0 & 0 & 0 & 0 & 0 & 0 \\
         1 & 0 & 0 & 0 & 0 & 1 & 0 \\
         2 & 0 & 0 & 0 & 1 & 0 & 0 \\
         3 & 0 & 0 & 0 & 1 & 1 & 1 \\
         4 & 0 & 0 & 1 & 0 & 0 & 0 \\
         5 & 0 & 0 & 1 & 0 & 1 & 1 \\
         6 & 0 & 0 & 1 & 1 & 0 & 1 \\
         7 & 0 & 0 & 1 & 1 & 1 & 0 \\
         8 & 0 & 1 & 0 & 0 & 0 & 1 \\
         9 & 0 & 1 & 0 & 0 & 1 & 1 \\
        10 & 0 & 1 & 0 & 1 & 0 & 0 \\
        11 & 0 & 1 & 0 & 1 & 1 & 0 \\
        12 & 0 & 1 & 1 & 0 & 0 & 0 \\
        13 & 0 & 1 & 1 & 0 & 1 & 1 \\
        14 & 0 & 1 & 1 & 1 & 0 & 0 \\
        15 & 0 & 1 & 1 & 1 & 1 & 0 \\
        16 & 1 & 0 & 0 & 0 & 0 & 1 \\
        17 & 1 & 0 & 0 & 0 & 1 & 1 \\
        18 & 1 & 0 & 0 & 1 & 0 & 1 \\
        19 & 1 & 0 & 0 & 1 & 1 & 1 \\
        20 & 1 & 0 & 1 & 0 & 0 & 0 \\
        21 & 1 & 0 & 1 & 0 & 1 & 0 \\
        22 & 1 & 0 & 1 & 1 & 0 & 1 \\
        23 & 1 & 0 & 1 & 1 & 1 & 0 \\
        24 & 1 & 1 & 0 & 0 & 0 & 1 \\
        25 & 1 & 1 & 0 & 0 & 1 & 1 \\
        26 & 1 & 1 & 0 & 1 & 0 & 0 \\
        27 & 1 & 1 & 0 & 1 & 1 & 0 \\
        28 & 1 & 1 & 1 & 0 & 0 & 0 \\
        29 & 1 & 1 & 1 & 0 & 1 & 0 \\
        30 & 1 & 1 & 1 & 1 & 0 & 0 \\
        31 & 1 & 1 & 1 & 1 & 1 & 0
    \end{tblr}
\end{table}

Переход всего клеточного автомата от поколения $t$ к поколению $t+1$
определяется выражением
\begin{equation}
    s^{t+1}_{i,j} = f\left(
        s^t_{i,j},
        s^t_{(i-1)\bmod n,j},
        s^t_{(i+1)\bmod n,j},
        s^t_{i,(j-1)\bmod n},
        s^t_{i,(j+1)\bmod n}
    \right).
    \label{eq:global-transition}
\end{equation}
Все значения $s^{t+1}_{i,j}$ в~\eqref{eq:global-transition} вычисляются только
по состоянию $C^t$. Поэтому переход является синхронным: новое состояние
одной клетки не влияет на вычисление остальных клеток в том же поколении.

Начальная конфигурация $C^0$ может быть задана явно либо сгенерирована
случайно. При случайной инициализации состояния клеток рассматриваются как
независимые бернуллиевские случайные величины с параметром $p$:
\begin{equation}
    \Pr\left(s^0_{i,j}=1\right)=p,
    \qquad
    \Pr\left(s^0_{i,j}=0\right)=1-p.
\end{equation}
После задания $C^0$ последовательность состояний
$C^0,C^1,C^2,\ldots$ однозначно определяется правилом~\eqref{eq:transition-rule}.

\subsection{Классы Вольфрама}

Для качественного описания динамики клеточных автоматов используется
классификация Вольфрама. В ней выделяются четыре класса поведения:
\begin{enumerate}
    \item эволюция почти всех начальных конфигураций приводит к однородному
        устойчивому состоянию;
    \item формируются устойчивые либо периодические структуры;
    \item возникает апериодическое хаотическое поведение, чувствительное к
        начальным условиям;
    \item появляются сложные локализованные структуры с длительным временем
        жизни, способные взаимодействовать между собой.
\end{enumerate}

Номер правила однозначно задаёт функцию перехода, однако сам по себе не
определяет класс поведения без исследования динамики. Для классификации
автомата необходимо рассматривать эволюцию различных начальных конфигураций,
наличие неподвижных и периодических состояний, распространение локальных
возмущений и изменение плотности живых клеток.

Для исследования правила подготовлены четыре начальные конфигурации,
представленные в табл.~\ref{tab:wolfram-presets}. Они демонстрируют режимы,
характерные для четырёх классов, но не означают принадлежность одного правила
одновременно к четырём классам: окончательная классификация выполняется по
совокупной динамике автомата.

\begin{table}[H]
    \centering
    \caption{Начальные конфигурации для исследования режимов поведения}
    \label{tab:wolfram-presets}
    \small
    \begin{tblr}{
        colspec = {Q[c,m,0.7] X[l,m,2.1] X[l,m,3.2]},
        row{1} = {font=\bfseries},
        hlines,
        vlines,
    }
        Режим & Начальная конфигурация & Наблюдаемое поведение \\
        I & В одном столбце жива каждая вторая клетка поля $30\times30$
          & Конфигурация остаётся неизменной \\
        II & Одна живая клетка в центре поля $25\times25$
           & После первого шага возникает структура с периодом 25 \\
        III & Случайное заполнение поля $25\times25$ с $p=0{,}5$ и seed 1
            & Повторение состояния не обнаружено за 5000 шагов \\
        IV & Крест из центральной клетки и четырёх соседей на поле $25\times25$
           & Длительная локализованная динамика; после 1154 шагов обнаружен цикл
             периода 100
    \end{tblr}
\end{table}
